\chapter[Proposed Methodology/Approach]{Proposed Methodology/Approach}{Proposed Methodology/Approach}\label{pm}

\rule{0.49\textwidth}{0.75pt} $_{\bigcirc}$ \rule{0.49\textwidth}{0.75pt}\\
\section{Problem Definition}
Federated Learning (FL) enables multiple healthcare institutions to collaboratively train a shared machine learning model without centralizing sensitive patient data. While this paradigm is well-aligned with clinical privacy requirements, practical deployments remain constrained by three interrelated challenges. First, communication overhead becomes substantial when high-dimensional gradients must be exchanged every round, limiting scalability and increasing latency in bandwidth-constrained settings. Second, Byzantine poisoning attacks---including gradient flipping, noise injection, model scaling, label-flip, and backdoor insertion---can significantly degrade model utility, and these attacks are particularly hard to detect under non-IID data distributions common across hospitals. Third, conventional FL provides no verifiable provenance of training activity: institutions cannot audit whether a global update was produced honestly, nor can they obtain immutable evidence of committee decisions and aggregation outcomes required for governance and regulatory traceability.

Therefore, there is a need for a secure, communication-efficient, and auditable FL workflow that can (i) reduce bandwidth cost while preserving similarity information for trust evaluation, (ii) filter malicious or low-quality updates under heterogeneous data, and (iii) record an immutable audit trail of each training round and its validation outcome.




\section{Objectives}
 \begin{enumerate}
  \item Make sure Privacy-Preserving Collaborative Training:
  \begin{enumerate}
    \item Enable distributed hospitals to train a shared diagnostic model without sharing raw patient data.
    \item Maintain privacy by confining training data locally at each institution while exchanging only model updates in a compressed and obfuscated form.
  \end{enumerate}
  
  \item Reduce Communication Burden:
  \begin{enumerate}
    \item Apply Random Projection-based gradient sketching to compress high-dimensional gradients into low-dimensional sketches while preserving directional similarity.
  \end{enumerate}
  
  \item Enhance Robustness Against Byzantine Clients:
  \begin{enumerate}
    \item Design an amnesic Data-Aware Verifier Selection (DAVS) mechanism that evaluates current-round gradient statistics to identify representative and trustworthy validators.
    \item Detect and exclude poisoned updates by jointly considering directional alignment (cosine similarity) and magnitude consistency.
  \end{enumerate}
  
  \item Provide Byzantine Fault-Tolerant Validation:
  \begin{enumerate}
    \item Integrate a PBFT-inspired consensus protocol among the selected committee to accept or reject global updates based on deterministic checks and quorum voting.
  \end{enumerate}
  
  \item Ensure Tamper-Proof Auditability:
  \begin{enumerate}
    \item Maintain a permissioned blockchain ledger that records per-round metadata (committee IDs, DAVS scores, PBFT votes, model hashes, and evaluation metrics) for immutable provenance.
  \end{enumerate}
  
  \item Validate Under Non-IID Medical Benchmarks:
  \begin{enumerate}
    \item Evaluate the end-to-end framework on non-IID medical imaging benchmarks (e.g., MedMNIST / PathMNIST) under diverse poisoning scenarios to demonstrate resilience and accuracy retention.
  \end{enumerate}
\end{enumerate}


\section{Scope}
DAVS targets privacy-sensitive healthcare settings where multiple institutions collaboratively train medical imaging models under non-IID data distributions. The scope covers the complete learning pipeline: local training, gradient sketching, committee formation via DAVS, PBFT-inspired validation, and blockchain-based logging of training provenance.

The framework focuses on communication efficiency (sketching high-dimensional updates into compact vectors) and Byzantine robustness (excluding poisoned updates and validating global updates through committee consensus). The blockchain layer is within a permissioned setting to provide auditability without the performance and openness constraints of public permissionless chains.

This scope does not attempt to solve every security threat in federated learning (e.g., full formal privacy guarantees or all adaptive collusion strategies). Instead, it establishes a unified, practical baseline that jointly addresses bandwidth reduction, poisoning resilience, and immutable traceability for regulated deployments.



\section{Proposed Approach}
The proposed approach implements DAVS as a secure federated learning workflow that couples gradient sketching, committee selection, consensus validation, and immutable logging.

In each training round, every hospital client performs local training for \(E\) epochs and computes its model update \(\Delta w_i = w_i^{local} - w^{global}\). To reduce bandwidth cost, the client compresses the high-dimensional update using a random projection matrix \(R\) to produce a compact sketch vector \(s_i = R \cdot \Delta w_i\). These sketches preserve similarity structure (enabling trust evaluation) while substantially reducing communication overhead and obfuscating raw gradient information.

After collecting sketches from all clients, the coordinator executes Data-Aware Verifier Selection (DAVS) to choose a small validator committee. DAVS is amnesic and operates only on current-round sketch statistics, ranking clients using pairwise cosine similarity (directional alignment) and magnitude deviation (consistency). The top-\(k\) clients are selected as the committee, while low-scoring outliers are treated as suspicious and excluded.

The selected committee then runs a PBFT-inspired protocol to validate the aggregated global update. Validators perform deterministic checks (e.g., no NaN/Inf values, bounded parameter changes, stable norms) and vote to accept or reject the update. If quorum is achieved, the update is committed; otherwise, the global model is retained for safety. Finally, a permissioned blockchain records the round outcome, model hash, committee membership, DAVS scores, PBFT votes, and evaluation metrics to provide tamper-proof auditability of the full training lifecycle.

\subsection{Proposed System Architecture}

\begin{figure}[h!]
    \centering
    \includegraphics[width=1\textwidth]{Chapter1/system_architecture.png} % Adjust the width as needed
    \caption{DAVS Framework Architecture}
    \label{f1} % Optional: for referencing the figure in your document
\end{figure}


The DAVS framework is designed as an end-to-end secure federated learning pipeline for healthcare institutions, emphasizing bandwidth efficiency, Byzantine robustness, and auditability. The architecture organizes each training round into modular stages: local training, gradient sketching, verifier (committee) selection, consensus-based validation, and blockchain logging. This modularity enables upgrades (e.g., alternate compression or validation checks) while preserving a consistent, verifiable workflow.

The system begins with the initialization of participating hospital clients and the coordinator. Each client maintains local medical datasets and trains a shared model without exporting raw data. In every round, clients compute updates and apply random projection compression to transmit low-dimensional sketches rather than full gradients. This reduces bandwidth significantly and supports scalability across distributed institutions.


Once sketches are collected, DAVS computes representativeness scores based on current-round alignment and magnitude consistency to form a validator committee. By not relying on historical reputation, the selection remains robust to adaptive attackers that attempt trust-building across rounds.

The selected committee then executes a PBFT-inspired consensus process to validate the aggregated model update before it becomes the new global model. The consensus layer ensures that even if some validators are faulty, the system can reject compromised updates and maintain global model integrity.

After validation, the audit layer appends a new block to a permissioned blockchain containing the round number, model hash, committee identifiers, DAVS scores, PBFT vote logs, and evaluation metrics. This yields a tamper-proof record of training provenance suitable for post-hoc analysis and compliance-oriented traceability in regulated healthcare environments.

Overall, the approach transforms a fragile, communication-heavy FL process into a verifiable and robust training pipeline. By combining sketching, amnesic verifier selection, PBFT-inspired validation, and blockchain logging, DAVS provides communication efficiency, Byzantine resilience, and immutable provenance in a single cohesive framework.




\rule{0.49\textwidth}{0.75pt} $_{\bigcirc}$ \rule{0.49\textwidth}{0.75pt}\\
